% ============================================================
%  ch03d  §3.4 Rankin--Selberg 动机的 Beilinson--Bloch--Kato 猜想
%  原文对应：ICM-arXiv.tex 1247--1321 行
%  编号：定理 3.4.1 (RS-Sym-Elliptic)；定理 3.4.2 (BBK for RS)
% ============================================================

\subsection{Rankin--Selberg 动机的 Beilinson--Bloch--Kato 猜想}\label{S: BBK}

设 $M$ 是数域 $F$ 上权为 $-1$ 的有理纯 Chow 动机。Beilinson--Bloch--Kato 猜想是 Birch--Swinnerton-Dyer 猜想的深远推广，它预言附于 $M$ 的若干代数、解析与上同调不变量之间的深刻联系：
\begin{itemize}
\item $M$ 的同调平凡圈的有理 Chow 群 $\CH(M)^0$；
\item $M$ 的 $L$-函数 $L(s,M)$；
\item $M$ 的 $\ell$-进实现 $H_{\ell}(M)$ 的 Bloch--Kato Selmer 群 $H^1_f(F, H_{\ell}(M))$。
\end{itemize}
Beilinson--Bloch 猜想预言
\[
\dim_{\bQ}\CH(M)^0 = \on{ord}_{s=0}L(s,M),
\]
即 $\CH(M)^0$ 的维数等于 $L$-函数在中心点的消没阶；而 Bloch--Kato 猜想预言
\[
\on{ord}_{s=0}L(s,M)= \dim_{\bQ_\ell} H^1_f(F, H_{\ell}(M)).
\]
此外，所谓的 $\ell$-进 Abel--Jacobi 映射
\[
\on{AJ}_\ell: \on{CH}(M)^0\otimes \bQ_\ell\to H^1_f(F, H_{\ell}(M))
\]
应当是同构。

该猜想目前看来完全无从下手。例如，对一般的动机，$L$-函数是否能亚纯延拓到整个复平面（从而 $L(s,M)$ 在 $s=0$ 处的消没阶才有意义）都还远未解决。（这一点会由数域上朗兰兹对应的 Galois 到自守方向推出。）尽管如此，已有许多工作在各种特殊情形下检验该猜想，多数是针对低秩的动机 $M$。在 \cite{LTXZZ-BBK} 中，我们对 Rankin--Selberg 动机验证了上述猜想的某些情形，而这类动机构成一列秩任意大的动机。

以下设 $F/F^+$ 是（非平凡的）CM 扩张，$F^+$ 全实。

\begin{thm}\label{RS-Sym-Elliptic}
设 $A_1,A_2$ 是 $F^+$ 上的两条椭圆曲线，并假设
\begin{enumerate}
\item\label{RS-Sym-Elliptic-1} $\End_{\overline F}A_i=\bZ$ 且 $\Hom_{\overline F}(A_1,A_2)=0$；
\item\label{RS-Sym-Elliptic-2} $\Sym^{n-1}A_1$ 与 $\Sym^nA_2$ 是模的；
\item\label{RS-Sym-Elliptic-3} 若 $n\geq 3$，则 $F^+\neq \bQ$。
\end{enumerate}
在这些假设下，若 $L(n, \Sym^{n-1}A_1\times \Sym^nA_2)\neq 0$，则对几乎所有的 $\ell$ 有
\[
\dim_{\bQ_\ell} H^1_f(F,  \Sym^{n-1}V_\ell(A_1)\otimes\Sym^nV_\ell(A_2)(1-n))=0.
\]
这里 $V_\ell(A_i)$ 照例记 $A_i$ 的有理 Tate 模。
\end{thm}

该定理其实是关于某些 Rankin--Selberg 自守表示所关联的 Galois 表示的 Bloch--Kato Selmer 群的一个更一般结果的特例，下面就来讨论它。

回忆一下，对 $\GL_{n}(\bA_F)$ 的一个不可约正则代数共轭自对偶尖（RACSDC）自守表示 $\Pi$，可以附上一族相容的 Galois 表示 $\rho_{\Pi,\la}:\Ga_F\to \GL_n(E_\la)$，其中 $E\subset \bC$ 是足够大的数域，$\la$ 是 $E$ 的一个素位（见 \cite{Chenevier-Harris-Galois-Rep}）。这样的 $E$ 称为 $\Pi$ 的强系数域；在下面考虑的情形中，它可以取为 $\Pi$ 的 Hecke 特征值生成的数域。

\begin{thm}\label{BBK for RS}
若 $n\geq 3$，设 $F^+\neq\bQ$。设 $\Pi_n$（相应地 $\Pi_{n+1}$）是 $\GL_{n}(\bA_F)$（相应地 $\GL_{n+1}(\bA_F)$）的 RACSDC 自守表示，无穷小特征平凡。设 $E\subseteq\bC$ 同时是 $\Pi_n$ 与 $\Pi_{n+1}$ 的强系数域。设 $\lambda$ 是 $E$ 的关于 $\Pi:=\Pi_0\times \Pi_1$ 的可容许素位\zy{原文此处把 $\Pi_n\times\Pi_{n+1}$ 写作 $\Pi_0\times\Pi_1$，按上下文应为前者；见勘误表。}。令 $\rho_{\Pi,\la}:=\rho_{\Pi_n,\la}\otimes_{E_\la} \rho_{\Pi_{n+1},\la}$。
\begin{enumerate}
\item\label{BBK for RS-1} 若 Rankin--Selberg $L$-值 $L(\frac{1}{2},\Pi)\neq 0$\footnote{这里用的是 $L$-函数的自守归一化。}，则 $H^1_f(F,\rho_{\Pi,\lambda}(n))=0$。
\item\label{BBK for RS-2} 若某个元素 $\Delta_\la\in H^1_f(F,\rho_{\Pi,\lambda}(n))$（下面会解释）非零，则 $H^1_f(F,\rho_{\Pi,\lambda}(n))$ 作为 $E_\la$-向量空间由 $\Delta_\la$ 生成。
\end{enumerate}
\end{thm}

定理中出现的可容许素位的概念包含一长串假设，其中一些相当技术性。本质上，它保证 Galois 表示 $\rho_{\Pi,\lambda}$ 有定义良好的 $\mO_{E,\la}$-格（下文仍记作 $\rho_{\Pi,\lambda}$），且模 $\la$ 约化表示足够大，并含某些特定的元素。（这也与定理 \ref{T:main theorem} 中 $V$-一般的概念有关。）所幸在一些有利情形下，可以证明除有限多个素位外都是可容许的。例如定理 \ref{RS-Sym-Elliptic} 所考虑的情形就是如此。纯自守框架下的另一个情形见 \cite[Thm. 1.1.7]{LTXZZ-BBK}。

该定理的证明用到几个不同的要素。情形 \eqref{BBK for RS-1} 的第一步，是借助整体 Gan--Gross--Prasad（GGP）猜想，把解析条件 $L(\frac{1}{2},\Pi)\neq 0$ 转化为更代数的条件。也就是说，GGP 猜想预言此时存在 $F$ 上的一对 Hermite 空间 $(V_n,V_{n+1})$，在 $\infty$ 处全正定，且 $V_{n+1}=V_n\oplus Fv$，其中 $(v,v)=1$；还存在酉群之积 $G=U(V_n)\times U(V_{n+1})$ 的一个温度尖自守表示 $\pi=\pi_n\times \pi_{n+1}$，使得周期积分
\[
[\Delta_H]: C^*_c(\Sh(G),E)[\pi]\to E
\]
不消失。这里 $H:=U(V_n)$ 对角地嵌入 $G$，从而诱导出相应志村簇（实际上是志村集）的嵌入 $\Delta_H: \Sh(H)\hookrightarrow \Sh(G)$，水平结构取适当（这里及以后都从记号中略去水平结构）。以 $[\Delta_H]$ 记 $\Sh(G)$ 中由 $\Sh(H)$ 给出的同调类，以 $C^*_c(\Sh(G),E)[\pi]$ 记 $\Sh(G)$ 的上同调（即函数空间）的 $\pi$-等型分支。

该猜想最初在 \cite{Zhang-GGP} 中，在某些局部限制下得到证明，而那些限制对算术应用来说过强。我们最近的工作借助迹公式研究中的一些新技术，把这些限制全部去掉了（\cite{BLZZ-GGP}）。

接下来的策略是：对每个 $m\geq 1$ 构造（无穷多个）上同调类 $\{\Theta_m^p\}_p\subset H^1(F, (\rho_{\Pi,\la}/\la^m)^*(1))$，其中 $p$ 是适当选取的素数，$(-)^*(1)$ 记 Pontryagin 对偶再扭以分圆特征；使得 $\Theta_m^p$ 与 Galois 表示 $\rho_{\Pi,\la}/\la^m$ 的 Selmer 类在 $p$ 处的局部 Tate 配对与上述周期积分相关。然后就可以用 Kolyvagin 型论证（在 \cite{Liu-HZ-Cycle, LTXZZ-BBK} 中作了放大），以 $\{\Theta_m^p\}$ 作为 Selmer 群的零化子，从而得出结论。

$\Theta_m^p$ 的构造用到志村簇的对角嵌入
\[
\Delta_{H'}: \Sh(H')\hookrightarrow \Sh(G'),
\]
其中 $H'\hookrightarrow G'$ 是 $H\subset G$ 的 prime-to-$p$ 平凡化（扩张纯）内形式（见定义 \ref{D: prime-to-p trivial}）。这些志村簇在 $p$ 处有抛孔水平结构，利用局部模型理论（\S\ref{S: local model}）可以证明它们的整模型在 $p$ 处是多重半稳定的，并算出其 mod $p$ 纤维上的邻近圈层。

综合运用许多要素——包括对 $\Sh(G')$ 在 $\overline{F}$ 上（整）上同调的理解、例 \ref{Tate cycle on U(1,2n)} 中的计算（特别是 \eqref{E:Hecke operator in Tpj} 与 \eqref{E: chern class on DL}）、以及 Taylor--Wiles 拼接方法 \cite{LTXZZ-Rigidity}——可以证明 $(\rho_{\Pi,\la}/\la^m)^*(1)$ 确实出现在 $\Sh(G')$ 的上同调中（即所谓的 $\Pi$ 的算术提高水平），并且对角圈 $\Delta_{H'}$ 在局部化到 $(\rho_{\Pi,\la}/\la^m)^*(1)$ 时确实给出所要的类 $\Theta_m^p$。需要说明的是，这与 \S\ref{S: Local Langlands} 与 \S\ref{geometric JL} 中的猜想是相容的：在 $\Loc_{\CG,p}\otimes \mO_E/\la^m$ 上，对应于 $c\on{-ind}_{K_p}^{G(\bQ_p)}(\mO_E/\la^m)$ 与 $c\on{-ind}_{K'_p}^{G'(\bQ_p)}(\mO_E/\la^m)$ 的凝聚层正应当以这种方式相关联。

情形 \eqref{BBK for RS-2} 中出现的类 $\Delta_\la$ 也可以解释。此时我们从一个志村簇的嵌入 $\Delta_H: \Sh(H)\hookrightarrow\Sh(G)$ 出发，其中 $G$ 是酉群之积，使得 $\Pi$ 下降为出现在 $\Sh_G$ 中间维上同调中的温度尖自守表示 $\pi$。于是圈 $\Delta_H$ 的 $\pi$-等型分支同调于零，我们令 $\Delta_\la=\on{AJ}_\la(\Delta_H[\pi])$。证明情形 \eqref{BBK for RS-2} 的策略，则是通过与前面类似的论证把它化归为情形 \eqref{BBK for RS-1}。
